
\documentclass{article}
\usepackage{colortbl}
\usepackage{makecell}
\usepackage{multirow}
\usepackage{supertabular}

\begin{document}

\newcounter{utterance}

\centering \large Interaction Transcript for game `wordle\_withclue', experiment `high\_frequency\_words\_clue\_no\_critic', episode 1 with reka{-}flash{-}3.1{-}t0.0.
\vspace{24pt}

{ \footnotesize  \setcounter{utterance}{1}
\setlength{\tabcolsep}{0pt}
\begin{supertabular}{c@{$\;$}|p{.15\linewidth}@{}p{.15\linewidth}p{.15\linewidth}p{.15\linewidth}p{.15\linewidth}p{.15\linewidth}}
    \# & \multicolumn{2}{c}{Player} && \multicolumn{2}{c}{Game Master} \\
    \hline

    \theutterance \stepcounter{utterance}  
    & & & \multicolumn{4}{p{0.6\linewidth}}{
        \cellcolor[rgb]{0.9,0.9,0.9}{
            \makecell[{{p{\linewidth}}}]{
                \texttt{\tiny{[P1$\langle$GM]}}
                \texttt{You are a language wizard who likes to guess words by using the given rules.} \\
\\ 
\texttt{Welcome to Wordle! You have six attempts to guess the target word, a valid English word of five lowercase letters (a{-}z). Please use the tags "explanation:" and "guess:" to provide a concise explanation for each guess.} \\
\\ 
\texttt{To help you make an informed guess, you will receive a clue for the word, such as} \\
\texttt{clue: snowy white.} \\
\\ 
\texttt{Here is an example guess based on the clue:} \\
\texttt{explanation: In the fairy tail Snow White, the girl is killed because she eats a poisoned apple. And the word apple has 5 letters.} \\
\texttt{guess: apple} \\
\\ 
\texttt{After each guess, your answer will be validated, and you will receive feedback indicating which letters are correct (green), which letters are correct but in the wrong position (yellow), and which letters are incorrect (red). This feedback can be useful in determining which letters to include or exclude in your next guess.} \\
\\ 
\texttt{For example, the feedback for "apple" might be:} \\
\texttt{guess\_feedback: a<yellow> p<yellow> p<green> l<yellow> e<red>} \\
\\ 
\texttt{The explanation should contain details about how both the guess\_feedback and clue work together in tandem to formulate a new guess.} \\
\\ 
\texttt{Do not generate any other text.} \\
\\ 
\texttt{Let's begin with your first guess.} \\
\\ 
\texttt{clue: comatose} \\
            }
        }
    }
     \\ \\

    \theutterance \stepcounter{utterance}  
    & \multicolumn{4}{p{0.6\linewidth}}{
        \cellcolor[rgb]{0.561,0.737,0.561}{
            \makecell[{{p{\linewidth}}}]{
                \texttt{\tiny{[P1$\rangle$GM]}}
                \texttt{<reasoning>} \\
\texttt{Let's tackle this Wordle with the clue "comatose". First, I need to understand what "comatose" means. It refers to a state of being deeply unconscious, right? So maybe the target word is related to that concept. But the target is a five{-}letter word. Let's break it down.} \\
\\ 
\texttt{The clue is "comatose", which is 7 letters. Maybe the target word is a shorter form or a part of that word. Let's see: "coma" is part of "comatose", but "coma" is four letters. The target needs to be five letters. Alternatively, maybe a synonym or related term. Words like "unconscious", "asleep", "stupor", but those are longer. Wait, maybe the target is a part of "comatose" but rearranged? Let's check letters: C, O, M, A, T, O, S, E. Five letters needed. Maybe "motes" (from "mote" but not sure). Alternatively, think of other words related to being unconscious. "Slump", "droop", "still"? Not sure.} \\
\\ 
\texttt{Alternatively, maybe the clue is a hint towards a different approach. "Comatose" could be a play on letters. Let's think of five{-}letter words that include letters from "comatose". For example, "motes" (m{-}o{-}t{-}e{-}s). Let's check if that's a valid word. Yes, "motes" means tiny particles. But how does that relate to "comatose"? Maybe not directly. Alternatively, "stoma" is part of "comatose" but "stoma" is five letters. Wait, "stoma" is an opening, like a mouth or a stomata (plural). But does "stoma" relate to being comatose? Not exactly. Maybe "motes" is a stretch.} \\
\\ 
\texttt{Alternatively, think of words that are associated with being in a coma. Words like "unconscious", "vegetative", but those are longer. Maybe "still" or "quiet"? Not sure. Alternatively, maybe the clue is a red herring, and the target is just a five{-}letter word that's part of "comatose". Let's list possible five{-}letter combinations from "comatose": "comat", "motes", "tomes", "somes", "a} \\
            }
        }
    }
    & & \\ \\

    \theutterance \stepcounter{utterance}  
    & & & \multicolumn{2}{p{0.3\linewidth}}{
        \cellcolor[rgb]{0.9,0.9,0.9}{
            \makecell[{{p{\linewidth}}}]{
                \texttt{\tiny{[GM$|$GM]}}
                \texttt{Error: The response should always start with the keyword 'explanation:'} \\
            }
        }
    }
    & & \\ \\

    \theutterance \stepcounter{utterance}  
    & & & \multicolumn{2}{p{0.3\linewidth}}{
        \cellcolor[rgb]{0.9,0.9,0.9}{
            \makecell[{{p{\linewidth}}}]{
                \texttt{\tiny{[GM$|$GM]}}
                \texttt{game\_result = ABORT} \\
            }
        }
    }
    & & \\ \\

\end{supertabular}
}

\end{document}
